\section*{अध्याय \ref{ch:g4:what-a-fire-needs} --- जलना: आग को क्या चाहिए}

\begin{solution}{exo:g4:what-a-fire-needs:1}
\omterm{def:g4:what-a-fire-needs:fuel}{ईंधन}, ऑक्सीजन (\omterm{def:g4:air-a-mixture-of-gases:air}{वायु} से) और ऊष्मा (आग शुरू करने के लिए)।
\end{solution}

\begin{solution}{exo:g4:what-a-fire-needs:2}
लकड़ी, मोमबत्ती का मोम, पेट्रोल और कागज़ \omterm{def:g4:what-a-fire-needs:fuel}{ईंधन} हैं। रेत, काँच और पानी
नहीं \omterm{def:g4:what-a-fire-needs:burning}{जलते}।
\end{solution}

\begin{solution}{exo:g4:what-a-fire-needs:3}
ऑक्सीजन: कंबल \omterm{def:g4:air-a-mixture-of-gases:air}{वायु} को आग तक पहुँचने से रोक देता है।
\end{solution}

\begin{solution}{exo:g4:what-a-fire-needs:4}
ऊष्मा: पानी \omterm{def:g4:what-a-fire-needs:burning}{जलती} लकड़ी को इतना ठंडा कर देता है कि वह अब
जल नहीं सकती।
\end{solution}

\begin{solution}{exo:g4:what-a-fire-needs:5}
इनमें से कोई भी तीन: राख, धुआँ, कार्बन डाइऑक्साइड, जलवाष्प।
\end{solution}

\begin{solution}{exo:g4:what-a-fire-needs:6}
फूँकने से गरम \omterm{def:g4:what-a-fire-needs:fuel}{ईंधन} तक ज़्यादा \omterm{def:g4:air-a-mixture-of-gases:air}{वायु}, यानी ज़्यादा ऑक्सीजन, पहुँचती है, और
वह फिर से \omterm{def:g4:what-a-fire-needs:burning}{जलने} लगता है।
\end{solution}

\begin{solution}{exo:g4:what-a-fire-needs:7}
पानी: \omterm{def:g4:what-a-fire-needs:burning}{जलता} मोम जलवाष्प बनाता है, जो ठंडे गिलास पर बूँदें
बन जाती है।
\end{solution}

\begin{solution}{exo:g4:what-a-fire-needs:8}
बाहर निकलना, धुएँ के नीचे झुककर और पीछे के दरवाज़े बंद करते हुए,
फिर बाहर से अग्निशमन दल को फ़ोन करना।
\end{solution}

\begin{solution}{exo:g4:what-a-fire-needs:9}
\omterm{def:g4:what-a-fire-needs:fuel}{ईंधन} वाली भुजा: गैस बंद होने पर आग के पास \omterm{def:g4:what-a-fire-needs:burning}{जलने} को कुछ नहीं
बचता।
\end{solution}

\begin{solution}{exo:g4:what-a-fire-needs:10}
\omterm{def:g4:what-a-fire-needs:fuel}{ईंधन}: पट्टी में पेड़ न होने से आग वहाँ पहुँचकर \omterm{def:g4:what-a-fire-needs:fuel}{ईंधन} से ख़ाली
रह जाती है।
\end{solution}

\begin{solution}{exo:g4:what-a-fire-needs:11}
नहीं। ज़्यादातर लकड़ी, \omterm{def:g4:air-a-mixture-of-gases:air}{वायु} की ऑक्सीजन के साथ मिलकर, ऐसे नए पदार्थ बन गई
है जो गैसें हैं --- कार्बन डाइऑक्साइड, जलवाष्प और धुआँ ---
और वे \omterm{def:g4:air-a-mixture-of-gases:air}{वायु} में चले गए हैं। केवल राख पीछे रहती है।
\end{solution}
